\chapter{Der Kanal, der zu viel puffert}\label{fall:kanal}

\begin{merksatz}
\boxtitel{Die Vorlage}
\textbf{Zusammenhang.} Übertragung, Go. Nimmt Messwerte von den Geräten entgegen und reicht sie an die Speicherung weiter.

\begin{center}
\begin{minipage}{0.9\textwidth}
\code{const puffergroesse = 500000}\\
\\
\code{eingang := make(chan Messwert, puffergroesse)}\\
\\
\code{go func() \{}\\
\code{\ \ \ \ for m := range eingang \{}\\
\code{\ \ \ \ \ \ \ \ speicher.Schreibe(m)}\\
\code{\ \ \ \ \}}\\
\code{\}()}
\end{minipage}
\end{center}

\begin{center}
\begin{minipage}{0.88\textwidth}
\kursiv{Der Puffer war eine Antwort auf ein echtes Problem: Vor zwei Jahren gingen bei Lastspitzen Messwerte verloren. Seit der Vergrößerung auf eine halbe Million ist das nicht mehr vorgekommen --- die Maßnahme hat gewirkt.}
\end{minipage}
\end{center}
\end{merksatz}

\begin{merksatz}
\boxtitel{Die Diagnose}
\begin{center}
{\small
\renewcommand{\arraystretch}{1.45}
\begin{tabular}{@{}P{2.4cm}P{2.0cm}P{2.0cm}P{2.2cm}@{}}
\toprule
\textbf{Füllstand} & \textbf{Speicher} & \textbf{Verzögerung} & \textbf{Verlust}\\
\midrule
Leer & $\mathbf{2}$ MB & $\mathbf{12}$ ms & $0$\\
$10\,\%$ & $18$ MB & $3{,}1$ s & $0$\\
$60\,\%$ & $86$ MB & \textcolor{akzent}{$\mathbf{19}$ \textbf{s}} & $0$\\
Voll & \textcolor{akzent}{$142$ MB} & \textcolor{akzent}{$\mathbf{31}$ \textbf{s}} & $0$\\
\bottomrule
\end{tabular}}
\end{center}

\begin{center}
\begin{minipage}{0.88\textwidth}
\kursiv{Kein Verlust in keiner Zeile --- die Maßnahme tut, was sie sollte. Und die dritte Spalte zeigt, was sie stattdessen anrichtet: Bei sechzig Prozent Füllstand erreicht ein Messwert die Klinik neunzehn Sekunden nach der Messung.}
\end{minipage}
\end{center}

\textbf{Was das klinisch bedeutet.}

\begin{center}
\renewcommand{\arraystretch}{1.45}
\begin{tabular}{@{}P{3.0cm}P{2.2cm}P{2.6cm}@{}}
\toprule
\textbf{Vorgang} & \textbf{Verzögerung erträglich?} & \\
\midrule
Verlaufsansicht & \textcolor{gruen}{\textbf{ja}} & $30$ s unerheblich\\
Auswertung nachts & \textcolor{gruen}{ja} & \\
\textbf{Kritische Grenzwertwarnung} & \textcolor{akzent}{\textbf{nein}} & \textcolor{akzent}{\textbf{$31$ s zu spät}}\\
\bottomrule
\end{tabular}
\end{center}

\begin{center}
\kursiv{Die dritte Zeile ist der Befund: Ein kritischer Kaliumwert löst die Warnung einunddreißig Sekunden später aus. Der Puffer hat den sichtbaren Verlust gegen eine unsichtbare Verzögerung getauscht --- und die Verzögerung erscheint in keiner Statistik.}
\end{center}
\end{merksatz}

\begin{technik}
\boxtitel{Das Urteil: umgebaut}
\begin{center}
\begin{minipage}{0.9\textwidth}
\code{dringend := make(chan Messwert)\ \ \ \ \ \ \ \ // ohne Puffer}\\
\code{normal\ \ := make(chan Messwert, 2000)}\\
\\
\code{go func() \{}\\
\code{\ \ \ \ for \{}\\
\code{\ \ \ \ \ \ \ \ select \{}\\
\code{\ \ \ \ \ \ \ \ case m := <-dringend:\ \ speicher.Schreibe(m)}\\
\code{\ \ \ \ \ \ \ \ default:}\\
\code{\ \ \ \ \ \ \ \ \ \ \ \ select \{}\\
\code{\ \ \ \ \ \ \ \ \ \ \ \ case m := <-dringend: speicher.Schreibe(m)}\\
\code{\ \ \ \ \ \ \ \ \ \ \ \ case m := <-normal:\ \ \ speicher.Schreibe(m)}\\
\code{\ \ \ \ \ \ \ \ \ \ \ \ \}}\\
\code{\ \ \ \ \ \ \ \ \}}\\
\code{\ \ \ \ \}}\\
\code{\}()}
\end{minipage}
\end{center}

\begin{center}
\begin{minipage}{0.86\textwidth}
\kursiv{Zwei Kanäle statt einem: kritische Werte ohne Puffer mit Vorrang, alles andere mit einem maßvollen Puffer. Die geschachtelte Auswahl erzwingt den Vorrang --- ein einfaches \code{select} über beide würde bei Bereitschaft würfeln.}
\end{minipage}
\end{center}
\end{technik}

\begin{merksatz}
\boxtitel{Die Begründung}
\begin{center}
{\small
\renewcommand{\arraystretch}{1.45}
\begin{tabular}{@{}P{2.8cm}P{2.0cm}P{2.2cm}P{2.0cm}@{}}
\toprule
& \textbf{Vorlage} & \textbf{Umgebaut} & \\
\midrule
Zeilen & $8$ & $\mathbf{18}$ & \\
Speicher & $142$ MB & $\mathbf{2{,}6}$ MB & $55\times$\\
\textbf{Verzögerung kritisch} & \textcolor{akzent}{$\mathbf{31}$ s} & \textcolor{gruen}{$\mathbf{14}$ \textbf{ms}} & \\
Verzögerung normal & $31$ s & $\mathbf{1{,}9}$ s & \\
Verlust bei Überlast & $0$ & \textcolor{mittel}{$\mathbf{0{,}3\,\%}$} & \textcolor{mittel}{gemeldet}\\
Aufwand & --- & $2$ PT & \\
\bottomrule
\end{tabular}}
\end{center}

\begin{center}
\begin{minipage}{0.88\textwidth}
\kursiv{Die vorletzte Zeile ist die unbequeme: Der Umbau bringt den Verlust zurück, den der große Puffer beseitigt hatte --- drei Promille der normalen Werte bei Lastspitzen. Dafür ist er gemeldet und betrifft nie einen kritischen Wert.}
\end{minipage}
\end{center}

\textbf{Warum das der bessere Handel ist.}

\begin{center}
\renewcommand{\arraystretch}{1.45}
\begin{tabular}{@{}P{2.6cm}P{2.6cm}P{2.8cm}@{}}
\toprule
& \textbf{Verlust} & \textbf{Verzögerung}\\
\midrule
Sichtbar & \textcolor{gruen}{\textbf{ja}} & \textcolor{akzent}{\textbf{nein}}\\
Zählbar & \textcolor{gruen}{ja} & \textcolor{mittel}{nur mit Messung}\\
Gerät kann puffern & \textcolor{gruen}{\textbf{ja}} & \textcolor{akzent}{nein}\\
Klinisch gefährlich & \textcolor{mittel}{begrenzt} & \textcolor{akzent}{\textbf{bei Warnungen}}\\
\bottomrule
\end{tabular}
\end{center}

\begin{center}
\kursiv{Die dritte Zeile entscheidet: Wenn der Puffer der Übertragung klein ist, greift der Gegendruck bis zum Gerät durch --- und das Gerät hat einen eigenen Puffer und kann entscheiden, was es verwirft. Band 22, Kapitel~5.}
\end{center}
\end{merksatz}

\begin{warnung}
\boxtitelr{Der Fehlschluss}
\begin{center}
\renewcommand{\arraystretch}{1.5}
\begin{tabular}{@{}P{3.2cm}P{2.2cm}P{2.4cm}@{}}
\toprule
\textbf{Naheliegend} & \textbf{Richtig?} & \\
\midrule
\enquote{Puffer sind schlecht} & \textcolor{akzent}{\textbf{nein}} & \textbf{$2000$ bleiben}\\
\enquote{Die frühere Änderung war falsch} & \textcolor{akzent}{\textbf{nein}} & \textcolor{akzent}{\textbf{sie war zu grob}}\\
\enquote{Kein Verlust ist immer besser} & \textcolor{akzent}{\textbf{nein}} & \textcolor{akzent}{das ist die Lehre}\\
\textbf{\enquote{Ein Puffer glättet, er ersetzt keine Kapazität}} & \textcolor{gruen}{\textbf{ja}} & \\
\bottomrule
\end{tabular}
\end{center}

\begin{center}
\begin{minipage}{0.88\textwidth}
\kursiv{Die zweite Zeile ist mir wichtig: Die Vergrößerung des Puffers war eine richtige Reaktion auf ein echtes Problem. Falsch war nur, dass sie alle Werte gleich behandelt hat --- die fachliche Unterscheidung fehlte, nicht die Maßnahme.}
\end{minipage}
\end{center}

\begin{center}
\kursiv{Die dritte Zeile ist die eigentliche Einsicht dieses Falls: Null Verlust bei einunddreißig Sekunden Verzögerung ist schlechter als drei Promille Verlust bei vierzehn Millisekunden --- wenn die betroffenen Werte Warnungen auslösen.}
\end{center}
\end{warnung}
\vspace{-4pt}
\noindent Klepp\-mann\zit{18} ordnet den Handel zwischen Ver\-lust und Ver\-zö\-ge\-rung ein; Nygard\zit{22} zeigt, warum groß\-zü\-gi\-ge Puffer den Zu\-sam\-men\-bruch nur ver\-schie\-ben.
